\section{Introduction}\label{sec:introduction}

A connected manifold is \emph{aspherical} if its universal cover is
contractible.  Its fundamental group then determines its homotopy type.
The topological Borel conjecture commonly asks whether every homotopy
equivalence between closed aspherical manifolds is homotopic to a
homeomorphism \cite{BartelsLueck2012,Lueck2010}.  The weaker
\emph{homeomorphism-existence formulation} asks only whether
homotopy-equivalent closed aspherical manifolds are homeomorphic.
Our main result disproves this formulation in dimension four.

\begin{theorem}\label{thm:pair}
There exist closed connected aspherical topological four-manifolds
$M$ and $N$ that are homotopy equivalent but not homeomorphic.
Their common fundamental group is word-hyperbolic.
\end{theorem}

The conclusion excludes every homeomorphism between the two manifolds,
without prescribing a homotopy class.  The construction also yields a
separate obstruction for a self-map of one manifold.

\begin{theorem}\label{thm:self-map}
There exist a closed connected aspherical topological four-manifold
$M$ and a homotopy equivalence $f\colon M\to M$ such that no
homeomorphism $M\to M$ is homotopic to $f$.
\end{theorem}

We prove this assertion separately and use the same $M$ in
Theorem~\ref{thm:pair}.  A nonhomeomorphic pair alone would not imply
the self-map assertion.

The homeomorphism-existence question already appears in Borel's letter
to Serre of May~2, 1953, phrased for compact manifolds that are
classifying spaces \cite{Borel1953}.  Rigidity holds classically in
dimensions at most two, and for closed orientable aspherical
three-manifolds by geometrization and three-dimensional rigidity
\cite[Theorem~0.7 and Section~5]{KreckLueck2009}.
In higher dimensions, Hsiang and Wall proved the homeomorphism-existence
statement for homotopy tori in dimensions at least five
\cite{HsiangWall1969}, and Farrell and Jones established it for closed
negatively curved manifolds in those dimensions
\cite{FarrellJones1989}.  Bartels and L\"uck proved the prescribed-class
formulation in dimensions at least five for a class of groups containing
word-hyperbolic groups and groups acting geometrically on
finite-dimensional CAT(0) spaces \cite[Theorem~A]{BartelsLueck2012}.

In dimension four, these rigidity theorems use further geometric input.  Bartels and
L\"uck's four-dimensional application assumes that the fundamental
group is good in the sense of Freedman
\cite[Proposition~0.3]{BartelsLueck2012}.  For closed aspherical
topological four-manifolds with torsion-free word-hyperbolic fundamental
group, Khan proves rigidity up to topological $s$-cobordism
\cite[Theorem~1.22 and Corollary~1.23]{Khan2012}; the product theorem for
five-dimensional topological $s$-cobordisms in Freedman--Quinn assumes
a good fundamental group \cite[Theorem~7.1A]{FreedmanQuinn1990}.
The word-hyperbolicity in Theorem~\ref{thm:pair} therefore lies within
the group-theoretic range of these rigidity results, while the
dimension is essential to their topological conclusion.  Recent
examples of Davis, Hayden, Huang, Ruberman, and Sunukjian are closed
aspherical smooth four-manifolds that are homeomorphic but not
diffeomorphic \cite[Theorem~1.1]{DHHRS2026}.  Their proof uses Davis
reflection and a characteristic-subgroup argument to force a hypothetical
diffeomorphism between these manifolds to lift to covers distinguished by
the genera of smoothly embedded homologically essential surfaces
\cite{DHHRS2026}.  Our conclusion concerns
homeomorphism in the topological category.

The proof uses two companion constructions.  The marked tensor
obstruction of \cite{MT} supplies plumbing blocks, marked grope bodies,
and algebraic potentials associated to three-dimensional cuts.  The
reflection and realization results of \cite{PD} turn the absence
of a marked filling into the absence of a closed manifold model for
an aspherical Poincar\'e complex.  We also use its universal-development and support arguments.
We state the imported results at their points of use and verify their
hypotheses for the chambers constructed here.  The principal inputs
are the marked algebraicity theorem \cite[Theorem~5.1]{MT}, the
controlled filling implication \cite[Sections~3--5]{PD}, and the
support construction \cite[Lemma~3.3]{PD}.

\subsection*{The construction}

Start with a finite collared Poincar\'e chamber $(P,S)$ of dimension
four, where $P$ has the homotopy type of a finite graph and its
fundamental group has an epimorphism to $\mathbb Z$.  Write $S_d$ for
the boundary in the connected cyclic cover of degree $d$.  A
\emph{marked filling} of $S_d$ is a compact four-manifold with this
identified boundary and the prescribed graph homotopy type, including
the boundary map.  The double cover $S_2$ has a marked filling, whereas
$S_d$ has none for odd $d\geq3$.

This parity comes from two colors of plumbing caps.  Each grope has
four terminal caps grouped into two red--blue pairs.  In the double
cover, choosing red caps at one vertex and blue caps at the other
makes the selected sheets disjoint and gives the filling.  For odd
$d\geq3$, a hypothetical filling of $S_d$ instead produces formal data
on a cycle of $d$ vertices with five red and five blue edges between
each neighboring pair.  The endpoint-assignment theorem would assign
every edge to one endpoint so that each vertex receives at most one
edge in some color.  Choose one such color at each vertex.  Adjacent
vertices cannot choose the same color, since together they would receive
at most two of the five edges of that color between them.  The choices
would alternate around the cycle, impossible when the cycle is odd.

We formulate the endpoint-assignment theorem for arbitrary finite
loopless multigraphs.
The cyclic-cover criterion works in every manifold dimension.  It
converts a double-cover model and eventual nonexistence of odd-cover
models into a self-homotopy equivalence not represented by any
homeomorphism.  The proof composes a lifted homeomorphism with a large
deck translation and imposes no finite-order condition on that
homeomorphism.  This last distinction separates the criterion from
finite-group Nielsen realization obstructions such as
\cite[Theorem~1.6]{BlockWeinberger2008}.

To construct the pair, use the reflection method of Davis
\cite{Davis1983,Davis2008} in the Poincar\'e-chamber form of
\cite{PD}.  A full reflection of $P$ has no manifold model in its
cyclic covers of odd degree $d\geq3$.  We quotient its elementary abelian
group of chamber labels by a binary relation space that acts freely
and forces every
simplicial symmetry preserving the relations to be trivial.  The
quotient is a finite aspherical complex $Q$ with the same odd-cover
obstruction.

There are two ways to replace the chambers of $Q_2$ by the filling of
$S_2$.  In one, every boundary identification is the same.  In the
other, it is twisted by the double-cover deck involution according to
a nonzero linear functional $\ell$ on the label group $\mathsf A$.
Both gluings are closed manifold models of $Q_2$.  They realize two
subgroups of outer automorphisms of the common group $\Pi$:
the horizontal subgroup $\mathsf A$ on $M$ and the graph subgroup
$\mathsf B_\ell$ on $N$.  These lie in
$\mathsf A\times\langle\delta\rangle\leq\Out(\Pi)$, where $\delta$ is
the cyclic involution.  The cyclic-cover criterion forbids $\delta$
on $M$ and at the same time proves Theorem~\ref{thm:self-map}.
The subgroups $\mathsf A$ and $\mathsf B_\ell$ are not conjugate in
$\Out(\Pi)$: they contain different numbers of elements admitting
finite-order lifts to $\Aut(\Pi)$.

The remaining step controls all possible homotopy classes of a
homeomorphism between the models.  We prove that $\Out(\Pi)$ is finite:
a uniform bound on square grids makes the support cubulation
hyperbolic, and Poincar\'e duality with the homological Shapiro lemma
excludes the virtually cyclic splittings forced by an infinite outer
automorphism group.  A direct rigidity argument for the right-angled
Coxeter group of the infinite lifted triangulation then gives
\[
 N_{\Out(\Pi)}(\mathsf A)=\mathsf A\times\langle\delta\rangle.
\]
The normalizer equality and the absence of $\delta$ make
$\mathsf A$ a Sylow $2$-subgroup of the outer automorphisms realized on
$M$.  A homeomorphism from $M$ to $N$ would transport
$\mathsf B_\ell$ into that realized group and force the forbidden
conjugacy.  This proves Theorem~\ref{thm:pair}.

Section~\ref{sec:cyclic} gives the two realization criteria.
Section~\ref{sec:reflection} states the chamber data and derives the
reflected odd-cover obstruction.  Sections~\ref{mod:section}
and~\ref{ext:section} construct the two models and identify their
outer actions.  Sections~\ref{hyp:section} and~\ref{nor:section}
prove the group-theoretic statements and the main theorem.
Finally, Sections~\ref{sec:chambers}--\ref{sec:formal} prove the
chamber data, including the general tensor theorem and the odd-cover
obstruction.

All manifolds are connected unless otherwise indicated.  Boundaries
have collars, and cornered constructions are rounded.  Homotopies
of markings are part of the data when boundary maps are prescribed.
We use singular cohomology unless a differential-form model is
specified, and write $\Gamma_j F$ for the lower central series of
a group $F$.
